
\subsection{Retrieval Tasks on Semi-structured Knowledge Bases}
\label{sec:star}
Our retrieval benchmark (Table~\ref{tab:benchmark_stats}) consists of three novel retrieval-based question-answering datasets, each comprising synthesized train/val/test sets with $9$k to $14$k queries in total and a high-quality human-generate query set.
% Specifically, the task is to retrieve items  over the SKBs given any input queries. 
The queries synthesize relational and textual knowledge, mirroring real-world queries in terms of natural-sounding property and flexible formats.

% Our retrieval benchmark, detailed in Table~\ref{tab:benchmark_stats}, comprises between $9$k to $14$k queries per dataset. Below, we outline the task on each dataset:

% Each query may have single answer or multiple answers, formulated as a subset of entities \revision{in the knowledge graph,}
% \jure{This is wrong. Entities are in the knowlede graph, Knowledge Base contains documents as well but those are not answers}
% which requires the model to accurately identify the ground truth entities out of potentially millions in response to each query.
% In Table~\ref{tab:benchmark_examples}, we demonstrate example queries from these benchmarks. 
\xhdr{\amazont} 
% This dataset comprises 9,100 synthesized customer-focused queries \hw{along with 81 human-generated queries for additional evaluation}. 
The task aims at product recommendation, with a notable 68\% of the synthesized queries yielding more than one ground truth answer. The dataset prioritizes customer-oriented criteria, highlighting textual elements such as product quality, functionality, and style. Moreover, it incorporate single-hop relational aspects (Appendix~\ref{app:amazon}) into the queries, including brand, category, and product connections (\eg complementary or substitute items). The queries are framed in conversation-like formats, enriching the context and enhancing the dataset's relevance to real-world scenarios.


% true positive: 86.6\%
% true negative: 96.7\%
\xhdr{\magt} 
% This dataset comprises 13,323 synthesized \hw{and 84 human-generated queries targeting at academic paper search.} 
Beyond the single-hop relational requirements in \amazont, \magh emphasizes the fusion between the textual requirements with multi-hop queries for precise academic paper search. For example, ``Are there any papers from King's College London'' highlights the metapath (\texttt{institution $\rightarrow$ author $\rightarrow$ paper}) on the relational structure. 
We designed three single-hop and four multi-hop relational query templates (Appendix~\ref{app:mag}). The textual aspects focus on the paper's topic, methodology, and contribution \etc
% true positive: 98.9\%

\xhdr{\primekgt} 
% This dataset, comprising 11,204 synthesized queries and \hw{109 human-generated queries across all ten entity types,} incorporates single-hop and multi-hop relational information similar to \magt. 
% For synthesized queries, we developed 28 multi-hop query templates, detailed in Appendix~\ref{app:prime}, to cover various relation types and ensure their practical relevance. 
% For instance, the query ``What is the drug that targets the genes or proteins expressed in <anatomy>?'' serves applications in precision medicine and pharmacogenomics, 
% aiding researchers and healthcare professionals in identifying drugs that act on genes or proteins associated with specific anatomical areas and enabling more targeted treatments. 
The task is to answer complex biomedicine inquiries. 
For synthesized queries, we developed 28 multi-hop query templates (Appendix~\ref{app:prime}) to cover various relation types and ensure their practical relevance. 
For example, the template ``What is the drug that targets genes or proteins in <anatomy>?'' aids precision medicine by identifying treatments targeted to specific anatomical areas.
For \texttt{drug, disease, gene/protein}, and \texttt{pathway} entities, the queries are a hybrid of relational and textual requirements. 
For entities such as \texttt{effect/phenotype}, the queries rely solely on relational data due to limited textual information. 
We exhibit three distinct user roles – medical scientist, doctor, and patient – for generating queries about drug and disease, which diversify the language to comprehensively evaluate the retrieval systems.



\begin{table}[t]
    \centering
    \caption{Statistics on the \benchmarkh benchmark datasets.}
    \resizebox{0.9\textwidth}{!}{
    \begin{tabular}{ll|rrrr}
        \toprule
        & & \multirow{2}{*}{\textbf{ \#queries}}  & \textbf{ \#queries w/} & \textbf{average} & \multirow{2}{*}{\textbf{train / val / test}}\\
        & & & \textbf{multiple answers}&  \textbf{\#answers} & \\
        \midrule
        \multirow{3}{*}{
        \makecell{Synthesized\\
        (Sec~\ref{sec:construction}, \ref{sec:analysis})}} & {\small \amazont} & 9,100 & 7,082  & 17.99& {0.65 / 0.17 / 0.18}\\
        & {\small \magt} &13,323 & 6,872 & 2.78 & {0.60 / 0.20 / 0.20}\\
        & {\small \primekgt} & 11,204 & 4,188 & 2.56 & {0.55 / 0.20 / 0.25} \\
        \midrule
        \multirow{3}{*}{\makecell{Human-generated\\
        (Sec~\ref{sec:human_generated})}} 
        & {\small \amazont} & 81 & 64 & 19.50 & \multirow{3}{*}{For testing only}\\
        & {\small \magt} & 84 &34 & 3.26 & \\
        & {\small \primekgt} & 98 & 41 & 2.77 & \\
        \bottomrule
        \end{tabular}
        }
    \vspace{-5pt}
    \label{tab:benchmark_stats}
\end{table}